\clearpage
\section{Разведочный анализ данных}
\label{sec:eda}

Разведочный анализ выполнен для компаний Y Combinator, имевших активный статус 13 июля 2023 года. Исходные характеристики взяты из выгрузки на эту дату, последующие статусы --- из выгрузки августовской версии 2026 года. Рассмотрены состав исследуемой группы, полнота записей и распределения характеристик. Источники данных опубликованы на Kaggle~\cite{kaggle-yc-dataset1,kaggle-yc-dataset2}.

\subsection{Формирование исследуемой совокупности}

В июльском файле из 8\,410 строк 4\,113 являются точными повторами. После их удаления остаются 4\,297 компаний с уникальными идентификаторами. На исходную дату 3\,110 из них были активны и вошли в исследуемую группу. Среди остальных 665 компаний были неактивны, 506 поглощены и 16 имели публичный статус. В поздней выгрузке найдены 3\,046 компаний исходной группы; 64 записи отсутствуют. Последним не присвоена отрицательная целевая метка, поскольку их статус в августовской версии 2026 года не установлен.

С февральской выгрузкой совпадают 4\,113 июльских идентификаторов. Только в феврале представлены 15 компаний, только в июле --- 184. Среди общих записей изменились 68 статусов, 357 значений размера команды и 62 года основания. В подсчёте учтены также переходы между заполненным и пустым значением. По двум выгрузкам нельзя отделить изменение характеристики от исправления карточки. Февральские значения не использованы для заполнения июльских пропусков.

Исходная группа ограничена компаниями, представленными в каталоге и сохранившими активный статус к июлю 2023 года. Поэтому её характеристики не описывают все компании, когда-либо принятые в Y Combinator, и тем более все предпринимательские инициативы.

\subsection{Полнота исходных характеристик}

На Рисунке~\ref{fig:eda-missing} приведены пропуски исходных характеристик. Год основания не указан у 356 компаний, что составляет 11,45\% исследуемой группы. Подробное описание отсутствует у 206 компаний, страна --- у 89, размер команды --- у 47, краткое описание --- у 43. Число основателей неизвестно у четырёх компаний после замены исходных нулей на пропуски.

\begin{figure}[!htbp]
    \centering
    \includegraphics[width=\textwidth]{eda/02_missing_values.pdf}
    \caption{Пропуски характеристик среди 3\,110 компаний Y Combinator, активных 13 июля 2023 года. Столбцы показывают число компаний без указанных сведений, подписи --- долю исходной группы. Пустые списки тегов учтены отдельно; данные августовской версии 2026 года для заполнения пропусков не использованы.}
    \label{fig:eda-missing}
\end{figure}

Пустой список тегов имеют 145 компаний, или 4,66\% группы. Для них число перечисленных меток равно нулю; сведения о деятельности компании могут при этом содержаться в описании. Длина описания и количество тегов характеризуют заполнение карточки, а не качество продукта.

Размер команды, число основателей, возраст, страна и краткое описание одновременно известны у 2\,632 компаний. Полное исключение строк с пропусками в этих пяти полях сократило бы исходную группу на 478 компаний, или 15,37\%. В описательных расчётах пропуски исключались отдельно для каждого показателя; в прогнозном анализе заполнение числовых значений выполнялось внутри обучающих частей.

\subsection{Последующие статусы и наблюдаемый показатель успеха}

В августовской версии 2026 года активный статус сохранили 2\,405 компаний исходной группы; 369 были неактивны, 265 поглощены, семь имели публичный статус. Распределение приведено на Рисунке~\ref{fig:eda-status}. Для 64 компаний поздний статус неизвестен. Зафиксированные статусы не содержат даты перехода между состояниями.

\begin{figure}[!htbp]
    \centering
    \includegraphics[width=\textwidth]{eda/03_followup_status.pdf}
    \caption{Статусы в августовской версии 2026 года у 3\,110 компаний Y Combinator, активных 13 июля 2023 года. Количество и доля каждой группы рассчитаны с учётом 64 компаний без известного позднего статуса. Положительная целевая метка объединяет поглощённые и публичные компании.}
    \label{fig:eda-status}
\end{figure}

Положительная целевая метка присвоена 272 компаниям, или 8,93\% от 3\,046 компаний с известным статусом. Остальные 2\,774 компании образуют отрицательный класс: на конечную дату они были активны либо неактивны, но не поглощены и не публичны. Эта метка не является общей оценкой экономического результата компании.

На поглощения приходятся 265 из 272 положительных исходов, или 97,43\%. Поэтому качество прогноза объединённого показателя в основном отражает распознавание поглощённых компаний. Семь публичных компаний представлены слишком малой группой для самостоятельной оценки качества прогноза этого исхода.

При двух крайних вариантах статуса 64 ненайденных компаний доля положительного исхода среди всех 3\,110 исходных компаний находится между $272/3\,110=8{,}75\%$ и $336/3\,110=10{,}80\%$. Эти границы отражают только неизвестные исходы и не учитывают возможные ошибки в записанных статусах.

Матрица переходов отдельно не строится. Все компании имели один исходный статус, поэтому её единственная содержательная строка совпала бы с распределением на Рисунке~\ref{fig:eda-status}. Кроме того, срезы не содержат дат изменения статуса и промежуточных состояний. Сопоставление показывает состояние компаний в двух точках времени, но не позволяет оценить скорость, последовательность или момент переходов.

\subsection{Годы набора и доступность последующих статусов}

Компании сгруппированы по году поступления в Y Combinator. На Рисунке~\ref{fig:eda-cohorts} сопоставлены численность групп и распределение поздних статусов внутри них. Ранние годы объединены, нестандартный код \texttt{IK12} выделен отдельно.

\begin{figure}[!htbp]
    \centering
    \includegraphics[width=\textwidth]{eda/04_cohort_composition.pdf}
    \caption{Численность и поздние статусы компаний Y Combinator по годам набора. Показаны 3\,110 компаний, активных 13 июля 2023 года; статусы взяты из августовской версии 2026 года. Слева указано число компаний, справа --- доли статусов внутри каждой группы, в том числе неизвестных. Ранние годы объединены; \texttt{IK12} оставлен отдельной группой.}
    \label{fig:eda-cohorts}
\end{figure}

К наборам 2021--2023 годов относятся 1\,680 компаний, или 54,02\% исходной группы. На 2018--2020 годы приходятся 845 компаний, на 2015--2017 годы --- 357, на стандартные наборы 2006--2014 годов --- 184. Ещё 44 компании имеют код \texttt{IK12}. Распределение отражает состав действующих компаний на июль 2023 года, а не число всех когда-либо принятых компаний: более ранние участники, утратившие активный статус до этой даты, отсутствуют.

Поздняя выгрузка не содержит 43 из 44 компаний группы \texttt{IK12}; доля ненайденных записей в ней составляет 97,73\%. Эта группа даёт 67,19\% всех неизвестных исходов. Среди остальных 3\,066 компаний не найдена 21, или 0,68\%. Следовательно, общая доля сопоставления 97,94\% скрывает сильную неоднородность по обозначению набора. Причина почти полного отсутствия \texttt{IK12} в позднем источнике по имеющимся данным не установлена.

У компаний с неизвестным поздним статусом чаще отсутствует год основания: 31 из 64 записей, или 48,44\%. Среди компаний с известным статусом этот показатель равен 325 из 3\,046, или 10,67\%. Поэтому выборка с известным исходом отличается от полной исходной группы не только доступностью целевой метки, но и полнотой характеристик.

Интервал между двумя срезами одинаков для всех компаний, но группы по году набора различаются возрастом и временем, прошедшим после поступления в акселератор. Наблюдаемые различия статусов не отделяют влияние года набора от этих характеристик.

\subsection{Распределения числовых характеристик}

На Рисунке~\ref{fig:eda-numerical} представлены размер команды, число основателей и приближённый возраст на июль 2023 года. Возраст рассчитан как разность календарных лет между исходным срезом и основанием компании.

\begin{figure}[!htbp]
    \centering
    \includegraphics[width=\textwidth]{eda/05_numerical_distributions.pdf}
    \caption{Распределения характеристик компаний Y Combinator, активных 13 июля 2023 года. Размер команды известен у 3\,063 компаний, число основателей --- у 3\,106, приближённый возраст --- у 2\,754. По вертикали показано число компаний. Для размера команды применена шкала $\ln(1+x)$ с подписями в количестве человек; столбцы показывают число компаний в интервалах, а не плотность.}
    \label{fig:eda-numerical}
\end{figure}

Медианный размер команды составляет девять человек, первый и третий квартили --- три и 25 человек. Среднее равно 55,59, максимум --- 7\,000. У 90\% компаний с заполненным полем команда насчитывает не более 80 человек. Из-за редких больших значений среднее заметно превышает медиану. Логарифмическая шкала позволяет показать небольшие команды без удаления 11 нулевых наблюдений и крайних значений; достоверность отдельных крупных значений по одному распределению не установлена.

Два основателя указаны у 1\,725 компаний, или 55,54\% записей с известным числом основателей. Один основатель указан у 692 компаний, три --- у 575. От четырёх до шести основателей имеют в сумме 114 компаний.

Медианный приближённый возраст равен трём календарным годам, первый и третий квартили --- двум и пяти годам; наблюдаемый диапазон составляет от нуля до 18 лет. Это возраст компаний, остававшихся активными на июль 2023 года, а не распределение срока существования всех компаний акселератора.

\subsection{Исходный размер команды в группах последующего статуса}

На Рисунке~\ref{fig:eda-boxplot} исходный размер команды сопоставлен с поздним статусом компании. Скрипичные диаграммы показывают распределение значений на логарифмической шкале; коробки и усы рассчитаны по исходному числу сотрудников. Чтобы редкие очень крупные команды не растягивали оценку плотности, значения выше 99-го процентиля в каждой группе показаны точками отдельно. При оценке квартилей эти значения сохранены.

\begin{figure}[!htbp]
    \centering
    \includegraphics[width=\textwidth]{eda/06_team_size_boxplot.pdf}
    \caption{Размер команды на 13 июля 2023 года по статусу компании в августовской версии 2026 года. Рассмотрены компании Y Combinator, активные на исходную дату; у каждой группы указано число записей с известным размером команды. Ширина скрипки показывает плотность на шкале $\ln(1+x)$ и нормирована отдельно для каждой группы. Значения выше 99-го процентиля исключены только из оценки плотности и показаны точками. Коробки и усы рассчитаны по всем исходным значениям: черта показывает медиану, коробка --- межквартильный интервал, усы --- крайние значения в пределах полутора межквартильных размахов. Семь публичных компаний показаны точками без оценки плотности. Подписи горизонтальной оси даны в количестве человек.}
    \label{fig:eda-boxplot}
\end{figure}

У компаний, сохранивших активный статус, медианный исходный размер команды равен десяти человек; у ставших неактивными --- пяти, у поглощённых --- восьми. Межквартильные интервалы составляют соответственно 4--30, 2--10 и 4--22,5 человека. Скрипичные диаграммы показывают концентрацию небольших команд во всех трёх группах и существенное перекрытие распределений. Поглощённые компании не выделяются более крупными исходными командами по сравнению с продолжающими работать.

В группе ставших неактивными у 76,0\% компаний с известным размером команды было не более десяти сотрудников. Среди сохранивших активный статус эта доля составляет 52,7\%, среди поглощённых --- 58,9\%. Обратная картина наблюдается в правой части распределения: команды численностью свыше 50 человек имели 3,9\% ставших неактивными компаний, 16,4\% сохранивших активный статус и 13,3\% поглощённых. Показатели рассчитываются функцией \texttt{boxplot} файла \href{https://github.com/Kiraashado/course_work_startup_success_predictors/blob/main/src/eda_analysis.py}{\texttt{src/eda\_analysis.py}} и сохраняются в поле \texttt{team\_by\_status} отчёта \path{eda_summary.json}. Таким образом, различия касаются не только медианы, но и доли крупных команд; сами по себе они не устанавливают влияния размера команды на последующий статус.

У семи компаний с последующим публичным статусом медианный размер команды на исходную дату равен 1\,050 человек, межквартильный интервал --- 175--2\,425. Эта малочисленная группа отличается по масштабу от остальных и делает положительный класс неоднородным. Различия на Рисунке не учитывают возраст и направление деятельности компаний.

\subsection{Зависимости между исходными числовыми признаками}

Матрица на Рисунке~\ref{fig:eda-correlation} содержит ранговые корреляции Спирмена между пятью исходными характеристиками. Коэффициент отражает согласованность порядка значений, а не линейное изменение одного признака при изменении другого. Каждая оценка рассчитана по компаниям с двумя известными значениями; число таких компаний указано в клетке.

\begin{figure}[!htbp]
    \centering
    \includegraphics[width=0.95\textwidth]{eda/07_correlation_matrix.pdf}
    \caption{Ранговые корреляции Спирмена между характеристиками 3\,110 компаний Y Combinator, активных 13 июля 2023 года. В каждой клетке указаны коэффициент и число компаний с обоими известными значениями. Насыщенность цвета соответствует модулю коэффициента, направление связи показывает его знак. Возраст рассчитан по разности календарных лет, длина описания измерена в символах; поздний статус в матрицу не включён.}
    \label{fig:eda-correlation}
\end{figure}

Наибольшая по модулю корреляция наблюдается между возрастом и размером команды: $\rho_S=0{,}58$ по 2\,716 компаниям. Более старшим компаниям чаще соответствует большая численность команды. Признаки связаны, но не взаимозаменяемы; сравнение команд между группами позднего статуса без учёта возраста может смешивать эти различия.

Корреляция возраста с числом тегов равна $\rho_S=-0{,}33$, размера команды с числом основателей --- $\rho_S=0{,}12$. Отрицательная связь возраста с числом тегов может отражать различия в заполнении карточек компаний разных лет; установить её причину по матрице нельзя.

Для каждой пары признаков из расчёта исключены только компании, у которых отсутствует хотя бы одно из двух значений. Поэтому число наблюдений различается между клетками и указано непосредственно под соответствующим коэффициентом. Расчёт выполняется функцией \texttt{correlation\_plot} в файле \href{https://github.com/Kiraashado/course_work_startup_success_predictors/blob/main/src/eda_analysis.py}{\texttt{src/eda\_analysis.py}}; числовые результаты сохраняются в файлах \path{spearman.csv} и \path{spearman_pair_counts.csv}. Матрица характеризует парные монотонные связи между признаками; она не показывает их дополнительный вклад в прогноз позднего статуса.

Целевая метка намеренно не включена в матрицу. Коэффициент между отдельным признаком и бинарным исходом показал бы только нескорректированную парную связь на всей выборке. Он не учитывал бы совместное действие признаков и не оценивал бы переносимость на не использованные при расчёте данные. Прогностическая полезность блоков характеристик проверяется на внешних разбиениях в разделе~\ref{sec:predictors}.

\subsection{Географический состав}

Географический состав исходной группы представлен на Рисунке~\ref{fig:eda-countries}. Для сравнения США со странами, имеющими несколько десятков наблюдений, использована логарифмическая шкала числа компаний.

\begin{figure}[!htbp]
    \centering
    \includegraphics[width=\textwidth]{eda/08_top_countries.pdf}
    \caption{Десять наиболее представленных стран среди 3\,110 компаний Y Combinator, активных 13 июля 2023 года. Положение точки показывает число компаний на логарифмической шкале, подпись --- количество и долю всей исходной группы. Страна взята из июльской выгрузки; у 89 компаний она отсутствует. Остальные страны на Рисунке не показаны.}
    \label{fig:eda-countries}
\end{figure}

США указаны у 2\,024 компаний, или 65,08\% исходной группы. Индия представлена 181 компанией, Канада --- 112, Великобритания --- 105. У следующих шести стран на Рисунке от 34 до 78 компаний. За пределами первой десятки находятся ещё 325 компаний; у 89 страна неизвестна. Доли рассчитаны по всем 3\,110 исходным компаниям, а не только по показанным странам.

На США приходится около двух третей исходной группы. Для большинства других стран число компаний невелико, а число положительных исходов ещё меньше. Показанное распределение характеризует присутствие стран в выборке, но не частоту поглощений или публичных статусов внутри них.

\subsection{Тематические теги исходных карточек}

На Рисунке~\ref{fig:eda-tags} показаны метки, перечисленные в июльских карточках компаний. Отраслевая классификация поздней выгрузки в исходные признаки не включена. Одна компания учитывается не более одного раза для каждого тега, но может присутствовать в нескольких строках Рисунка.

\begin{figure}[!htbp]
    \centering
    \includegraphics[width=\textwidth]{eda/09_top_tags.pdf}
    \caption{Пятнадцать наиболее частых меток у 3\,110 компаний Y Combinator, активных 13 июля 2023 года. Показаны число компаний с каждой меткой и доля всей исходной группы. Одна компания может иметь несколько меток, поэтому доли не суммируются до 100\%. Обозначения \texttt{AI} и \texttt{Artificial Intelligence} сохранены раздельно; у 145 компаний список меток пуст.}
    \label{fig:eda-tags}
\end{figure}

Чаще всего встречается метка \texttt{SaaS}, обозначающая программное обеспечение как услугу: её имеют 819 компаний, или 26,33\% группы. Метка \texttt{B2B} для решений, адресованных организациям, указана у 710 компаний, метка \texttt{Fintech} для финансовых технологий --- у 559, инструменты для разработчиков --- у 338. Эти обозначения описывают разные стороны деятельности и не образуют взаимоисключающую отраслевую классификацию.

Несколько меток относятся к искусственному интеллекту: \emph{Artificial Intelligence} встречается у 269 компаний, \texttt{AI} --- у 130, \emph{Generative AI} --- у 114; метка машинного обучения указана у 173 компаний. Эти группы могут пересекаться, поэтому сумма чисел не показывает количество компаний, работающих с искусственным интеллектом. Исходные обозначения на Рисунке не объединены.

В прогнозном эксперименте исходные тематические метки кодируются отдельно без объединения в смысловые группы. Словарь формируется только по обучающей части каждого разбиения.
\clearpage
